\documentclass[%
    a4paper,
    oneside,
    punct=kaiming,
    zihao=-4,
]{ctexbook}

\providecommand*{\laauxtoro}{the \textsf{mwe}
    (Minimal Working Example) package}
\providecommand*{\lailustrajxo}{example-image-a4.pdf}

\usepackage[margin=2.5cm]{geometry}
\pagestyle{empty}

\IfFontExistsTF{FiraSans-Regular.otf}{%
    \IfFontExistsTF{FiraSans-SemiBold.otf}{%
        \IfFontExistsTF{FiraSans-Italic.otf}{%
            \IfFontExistsTF{FiraSans-SemiBoldItalic.otf}{%
                \setmainfont{FiraSans}[
                    Path,
                    UprightFont = *-Regular.otf,
                    BoldFont = *-SemiBold.otf,
                    ItalicFont = *-Italic.otf,
                    BoldItalicFont = *-SemiBoldItalic.otf
                ]
                \setsansfont{FiraSans}[
                    Path,
                    UprightFont = *-Regular.otf,
                    BoldFont = *-SemiBold.otf,
                    ItalicFont = *-Italic.otf,
                    BoldItalicFont = *-SemiBoldItalic.otf
                ]
            }{}
        }{}
    }{}
}{}

\IfFontExistsTF{Sarasa Mono SC}{%
    \IfFontExistsTF{Sarasa Mono SC Italic}{%
        \IfFontExistsTF{Sarasa Mono SC SemiBold}{%
            \IfFontExistsTF{Sarasa Mono SC SemiBold Italic}{%
                \PassOptionsToClass{fontset=none}{ctexbook}
                \setmonofont{Sarasa Mono SC}[
                    UprightFont = *,
                    BoldFont = * SemiBold,
                    ItalicFont = * Italic,
                    BoldItalicFont = * SemiBold Italic
                ]
                \setCJKmainfont{Sarasa Mono SC}[
                    UprightFont = *,
                    BoldFont = * SemiBold,
                    ItalicFont = * Italic,
                    BoldItalicFont = * SemiBold Italic
                ]
                \setCJKsansfont{Sarasa Mono SC}[
                    UprightFont = *,
                    BoldFont = * SemiBold,
                    ItalicFont = * Italic,
                    BoldItalicFont = * SemiBold Italic
                ]
                \setCJKmonofont{Sarasa Mono SC}[
                    UprightFont = *,
                    BoldFont = * SemiBold,
                    ItalicFont = * Italic,
                    BoldItalicFont = * SemiBold Italic
                ]
            }{}
        }{}
    }{}
}{
    \IfFontExistsTF{SarasaMonoSC-Regular.ttf}{%
        \IfFontExistsTF{SarasaMonoSC-Italic.ttf}{%
            \IfFontExistsTF{SarasaMonoSC-SemiBold.ttf}{%
                \IfFontExistsTF{SarasaMonoSC-SemiBoldItalic.ttf}{%
                    \PassOptionsToClass{fontset=none}{ctexbook}
                    \setmonofont{SarasaMonoSC}[
                        Path,
                        UprightFont = *-Regular.ttf,
                        BoldFont = *-SemiBold.ttf,
                        ItalicFont = *-Italic.ttf,
                        BoldItalicFont = *-SemiBoldItalic.ttf
                    ]
                    \setCJKmainfont{SarasaMonoSC}[
                        Path,
                        UprightFont = *-Regular.ttf,
                        BoldFont = *-SemiBold.ttf,
                        ItalicFont = *-Italic.ttf,
                        BoldItalicFont = *-SemiBoldItalic.ttf
                    ]
                    \setCJKsansfont{SarasaMonoSC}[
                        Path,
                        UprightFont = *-Regular.ttf,
                        BoldFont = *-SemiBold.ttf,
                        ItalicFont = *-Italic.ttf,
                        BoldItalicFont = *-SemiBoldItalic.ttf
                    ]
                    \setCJKmonofont{SarasaMonoSC}[
                        Path,
                        UprightFont = *-Regular.ttf,
                        BoldFont = *-SemiBold.ttf,
                        ItalicFont = *-Italic.ttf,
                        BoldItalicFont = *-SemiBoldItalic.ttf
                    ]
                }{}
            }{}
        }{}
    }{}
}

\usepackage[
    math-style=ISO,
    bold-style=ISO,
    mathrm=sym,
    mathbf=sym,
]{unicode-math}
\IfFontExistsTF{FiraMath-Regular.otf}{%
    \setmathfont{FiraMath-Regular.otf}
}{}

\makeatletter
\providecommand*{\midsloppy}{%
    \tolerance 5000%
    \hbadness 4000%
    \emergencystretch 1.5em%
    \hfuzz .1\p@%
    \vfuzz\hfuzz}
\makeatother

\renewcommand*{\familydefault}{\sfdefault}

\AtBeginDocument{%
    \frenchspacing
    \midsloppy
    \renewcommand*{\UrlFont}{\ttfamily}
}

\usepackage[breaklinks,hidelinks,unicode]{hyperref}
\usepackage{xurl}

\usepackage{caption}
\usepackage{subcaption}
\usepackage{qrcode}
\qrset{link,height=2.5cm}
\usepackage{tikz}
\usepackage{pdfpages}

\begin{document}

\includepdf{\lailustrajxo}
\clearpage

\begin{align*}
         &
    \det \begin{bmatrix}
             a_{1,1} & a_{1,2} & a_{1,3} \\
             a_{2,1} & a_{2,2} & a_{2,3} \\
             a_{3,1} & a_{3,2} & a_{3,3} \\
         \end{bmatrix}
    \\
    = {} &
    a_{1,1} a_{2,2} a_{3,3}
    + a_{1,2} a_{2,3} a_{3,1}
    + a_{1,3} a_{2,1} a_{3,2}
    - a_{1,1} a_{2,3} a_{3,2}
    - a_{1,2} a_{2,1} a_{3,3}
    - a_{1,3} a_{2,2} a_{3,1}
\end{align*}

\begin{figure*}[h!]
    % Postscript on 2024-09-29:
    % The following figure is a derivative of
    % the work by "越来越少"
    % and the work by "Sagittarius Rover".
    % For more information, visit
    % https://ask.latexstudio.net/ask/question/17406.html
    % I thank them very much.
    \centering
    % \begin{tikzpicture}
    \begin{tikzpicture}[scale=1.5]
        \colorlet{verda}{green!70!black}
        \colorlet{rugxa}{red!80}

        \tikzset{
            temp/.cd,
            s opt/.store in=\solidlineopt, s opt={draw=verda},
            s path/.store in=\solidlinepath, s path=,
            s end opt/.store in=\solidlineendopt, s end opt={right},
            s end text/.store in=\solidlineendtext, s end text=,
            d opt/.store in=\dashedlineopt, d opt={draw=rugxa,dashed},
            d path/.store in=\dashedlinepath, d path=,
            d end opt/.store in=\dashedlineendopt, d end opt={left},
            d end text/.store in=\dashedlineendtext, d end text=,
        }
        \def\linesdraw#1{
            \begingroup
            \tikzset{temp/.cd,#1}
            \edef\tempcmd{
                \noexpand\begin{scope}
                    \noexpand\draw[\solidlineopt]\solidlinepath node[\solidlineendopt]{\solidlineendtext};
                    \noexpand\end{scope}
                \noexpand\begin{scope}
                    \noexpand\draw[yscale=-1,rotate=-90,\dashedlineopt]
                    \dashedlinepath
                    node[\dashedlineendopt]{\dashedlineendtext};
                    \noexpand\end{scope}
            }
            \tempcmd
            \endgroup
        }

        \begin{scope}[rotate=-135]
            \foreach \i/\j/\k in {
            {(4,4)--(4,0)arc(0:-180:1.5)--(1,5)}/{\textcolor{verda}{$+a_{1,2}a_{2,3}a_{3,1}$}}/{\textcolor{rugxa}{$-a_{1,2}a_{2,1}a_{3,3}$}},
            {(3,3)--(3,0)arc(0:-180:1.5)--(0,5)}/{\textcolor{verda}{$+a_{1,3}a_{2,1}a_{3,2}$}}/{\textcolor{rugxa}{$-a_{1,1}a_{2,3}a_{3,2}$}},
            {(2,0)--(2,5)}/{\textcolor{verda}{$+a_{1,1}a_{2,2}a_{3,3}$}}/{\textcolor{rugxa}{$-a_{1,3}a_{2,2}a_{3,1}$}}
            }{
            \linesdraw{
                %  s opt=,
                s path={\i},
                %  s end opt={},
                s end text={\j},
                %  d opt=,
                d path={\i},
                %  d end opt={},
                d end text={\k},
            }}
        \end{scope}

        \begin{scope}[yscale={-sqrt(2)},xscale={sqrt(2)},shift={(-2,0)}]
            \foreach \i in {1,2,3}{
                    \foreach \j in {1,2,3}{
                            \node [fill=white] at(\i,\j){\(a_{\j,\i}\)};
                        }}
        \end{scope}
    \end{tikzpicture}

    \caption*{%
        \centering%
        Memorhelpilo de la determinanto de \(3 \times 3\)-matrico\\
        A mnemonic device of the determinant of a \(3 \times 3\) matrix\\
        \(3 \times 3\) 阵的行列式的记忆方法}
\end{figure*}

\begin{figure*}[h!]
    \centering
    \begin{subfigure}{0.4\textwidth}
        \centering
        \qrcode{https://github.com/septsea/det}
        \caption*{\url{https://github.com/septsea/det}}
    \end{subfigure}
    \begin{subfigure}{0.4\textwidth}
        \centering
        \qrcode{https://gitee.com/septsea/det}
        \caption*{\url{https://gitee.com/septsea/det}}
    \end{subfigure}

    \caption*{%
        \centering%
        Legu \textit{Enkonduko al determinantoj} por pli multaj informoj.\\
        Read \textit{An introduction to determinants} for more information.\\
        阅读\,《行列式入门》\,以了解更多.}
\end{figure*}

\vfill

\begingroup

\setlength{\parindent}{0pt}

\footnotesize

\texttt{©} Determinantulo.
All rights reserved.

The illustration on the back of the paper was drawn by
\laauxtoro%
.

% \textit{An introduction to determinants} is available under
% the Zero-Clause BSD License.
% % the Zero-Clause BSD License
% % \url{https://opensource.org/license/0bsd}.

\endgroup

\end{document}
